% ═══════════════════════════════════════════════════════════════════
% §5  Tensor Metadata
% ═══════════════════════════════════════════════════════════════════

\section{Tensor Metadata}
\label{sec:tensormeta}

\subsection{TensorMeta Structure}

\begin{definition}[TensorMeta]
\label{def:tensormeta}
A compact metadata record for one tensor, capturing shape, strides,
dtype, device, and data pointer without any actual tensor data.
Exactly $\bytes{144}$.
\end{definition}

\begin{table}[h]
\centering
\begin{tabular}{r r l l l}
\toprule
\textbf{Offset} & \textbf{Size} & \textbf{Field} & \textbf{Type} & \textbf{Default} \\
\midrule
  0 & 64B & \field{sizes[8]}      & \code{int64\_t[8]}  & $\{0\}$ \\
 64 & 64B & \field{strides[8]}    & \code{int64\_t[8]}  & $\{0\}$ \\
128 &  8B & \field{data\_ptr}     & \code{void*}        & \code{nullptr} \\
136 &  1B & \field{ndim}          & \code{uint8\_t}     & $0$ \\
137 &  1B & \field{dtype}         & \struct{ScalarType}  & \code{Undefined} \\
138 &  1B & \field{device\_type}   & \struct{DeviceType} & \code{CPU} \\
139 &  1B & \field{device\_idx}   & \code{int8\_t}      & $-1$ \\
140 &  1B & \field{layout}        & \struct{Layout}     & \code{Strided} \\
141 &  3B & \field{pad[3]}        & \code{uint8\_t[3]}  & $\{0\}$ \\
\midrule
    & 144B & \multicolumn{3}{l}{\textit{Total (verified by \code{static\_assert})}} \\
\bottomrule
\end{tabular}
\caption{TensorMeta layout.  The \field{sizes} and \field{strides} arrays
support up to 8 dimensions, covering 99.9\% of real tensors.}
\end{table}

\begin{invariant}[Dimension bound]
$0 \le \field{ndim} \le 8$.  Tensors with more than 8 dimensions
require indirection at a higher level (not inside this struct).
\end{invariant}

\subsection{Storage Size Computation}
\label{sec:storage-nbytes}

\begin{definition}[\code{compute\_storage\_nbytes}]
\label{def:storage-nbytes}
Given a \struct{TensorMeta} $m$, the storage size in bytes is:
\[
  \operatorname{nbytes}(m) = \bigl(\max\_\text{off} - \min\_\text{off} + 1\bigr)
    \times \operatorname{element\_size}(m.\field{dtype})
\]
where the offset extremes are computed as:
\begin{align*}
  \max\_\text{off} &= \sum_{d=0}^{m.\field{ndim}-1}
    \max\!\bigl(0,\; (m.\field{sizes}[d] - 1) \times m.\field{strides}[d]\bigr) \\[4pt]
  \min\_\text{off} &= \sum_{d=0}^{m.\field{ndim}-1}
    \min\!\bigl(0,\; (m.\field{sizes}[d] - 1) \times m.\field{strides}[d]\bigr)
\end{align*}
\end{definition}

\begin{algorithm_spec}[\code{compute\_storage\_nbytes}]
\label{alg:storage-nbytes}
\begin{enumerate}
  \item If $\field{ndim} = 0$: return $\operatorname{element\_size}(\field{dtype})$ (scalar tensor).
  \item Initialize $\max\_\text{off} \leftarrow 0$, $\min\_\text{off} \leftarrow 0$.
  \item For each dimension $d \in [0, \field{ndim})$:
    \begin{enumerate}
      \item If $\field{sizes}[d] = 0$: return $0$ (zero-size tensor).
      \item $\text{extent} \leftarrow (\field{sizes}[d] - 1) \times \field{strides}[d]$.
      \item If $\text{extent} > 0$: $\max\_\text{off} \mathrel{+}= \text{extent}$.
      \item Else: $\min\_\text{off} \mathrel{+}= \text{extent}$.
    \end{enumerate}
  \item Return $(\max\_\text{off} - \min\_\text{off} + 1) \times
    \operatorname{element\_size}(\field{dtype})$.
\end{enumerate}
\end{algorithm_spec}

\begin{remark}[Negative strides]
Negative strides arise from \code{torch.flip} and \code{as\_strided}.
For example, a reversed vector of length $n$ with stride $-1$ has
$\min\_\text{off} = -(n-1)$, $\max\_\text{off} = 0$, giving
$\operatorname{nbytes} = n \times \operatorname{element\_size}$.
The tracking of both extremes is essential for correctness.
\end{remark}

\begin{property}[Consistency with PyTorch]
For contiguous tensors: $\operatorname{nbytes}(m)
= \bigl(\prod_{d} m.\field{sizes}[d]\bigr) \times \operatorname{element\_size}(m.\field{dtype})$.
The general formula handles non-contiguous layouts
where $\text{stride} \neq \text{size product of trailing dimensions}$.
\end{property}

\subsection{View and Alias Semantics}

\begin{definition}[View tensor]
A view tensor shares the same \field{data\_ptr} as its base tensor
but has different \field{sizes} and/or \field{strides}.
Two \struct{TensorMeta} entries are \emph{aliases} if and only if
their \field{data\_ptr} fields are equal and at least one of
\field{sizes} or \field{strides} differs.
\end{definition}

Alias detection is performed by the PtrMap during graph construction (\S\ref{sec:graph-construction}).
Same \field{data\_ptr} from different ops $\to$ an \code{ALIAS} edge.

\subsection{External Tensor Classification}

\begin{definition}[External tensor]
\label{def:external-tensor}
A tensor whose storage is not produced by any recorded op in the
current iteration: model parameters, data loader outputs, and
tensors created before recording began.  Identified by
$\field{op\_index} = \text{OpIndex::none()} = 2^{32}-1$ in the PtrMap.
\end{definition}

External tensors are excluded from the memory pool
(they keep their existing allocations) and marked with
$\field{is\_external} = \text{true}$ in their \struct{TensorSlot}.

\subsection{TensorSlot Structure}
\label{sec:tensorslot}

\begin{definition}[TensorSlot]
\label{def:tensorslot}
One unique storage in a recorded iteration, tracking its live range
and assigned memory pool offset.  Exactly $\bytes{40}$.
\end{definition}

\begin{table}[h]
\centering
\begin{tabular}{r r l l l}
\toprule
\textbf{Offset} & \textbf{Size} & \textbf{Field} & \textbf{Type} & \textbf{Default} \\
\midrule
 0 &  8B & \field{offset\_bytes}  & \code{uint64\_t}    & $0$ \\
\midrule
\multicolumn{5}{l}{\textit{--- Bulk-copyable region (24B, matches SlotInfo) ---}} \\
 8 &  8B & \field{nbytes}        & \code{uint64\_t}    & $0$ \\
16 &  4B & \field{birth\_op}     & \struct{OpIndex}    & none \\
20 &  4B & \field{death\_op}     & \struct{OpIndex}    & none \\
24 &  1B & \field{dtype}         & \struct{ScalarType}  & \code{Undefined} \\
25 &  1B & \field{device\_type}   & \struct{DeviceType} & \code{CPU} \\
26 &  1B & \field{device\_idx}   & \code{int8\_t}      & $-1$ \\
27 &  1B & \field{layout}        & \struct{Layout}     & \code{Strided} \\
28 &  1B & \field{is\_external}  & \code{bool}         & \code{false} \\
29 &  3B & \field{pad[3]}        &                     & $\{0\}$ \\
\midrule
32 &  4B & \field{slot\_id}      & \struct{SlotId}     & none \\
36 &  4B & \field{pad2[4]}       &                     & $\{0\}$ \\
\midrule
   & 40B & \multicolumn{3}{l}{\textit{Total (verified by \code{static\_assert})}} \\
\bottomrule
\end{tabular}
\caption{TensorSlot layout.  The bulk-copyable region at bytes 8--31
matches the \struct{SlotInfo} scratch buffer layout, enabling a single
$\bytes{24}$ \code{memcpy} during graph construction.}
\end{table}

\begin{definition}[Lifetime interval]
\label{def:lifetime}
A tensor slot $s$ is \emph{alive} at op $i$ when:
\[
  s.\field{birth\_op}.\text{raw()} \le i \le s.\field{death\_op}.\text{raw()}
\]
The birth is the producing op index; the death is the last consuming op index.
External tensors have $\field{birth\_op} = \text{none}$ (alive for the
entire iteration).
\end{definition}
